
\documentclass{article}
\usepackage{colortbl}
\usepackage{makecell}
\usepackage{multirow}
\usepackage{supertabular}

\begin{document}

\newcounter{utterance}

\centering \large Interaction Transcript for game `eqbench', experiment `eq\_bench', episode 79 with reflexion{-}Qwen3.5{-}27B.
\vspace{24pt}

{ \footnotesize  \setcounter{utterance}{1}
\setlength{\tabcolsep}{0pt}
\begin{supertabular}{c@{$\;$}|p{.15\linewidth}@{}p{.15\linewidth}p{.15\linewidth}p{.15\linewidth}p{.15\linewidth}p{.15\linewidth}}
    \# & \multicolumn{2}{c}{Player} && \multicolumn{2}{c}{Game Master} \\
    \hline

    \theutterance \stepcounter{utterance}  
    & & & \multicolumn{4}{p{0.6\linewidth}}{
        \cellcolor[rgb]{0.9,0.9,0.9}{
            \makecell[{{p{\linewidth}}}]{
                \texttt{\tiny{[P1$\langle$GM]}}
                \\ 
\texttt{Your task is to predict the likely emotional responses of a character in this dialogue:} \\
\\ 
\texttt{Bennett: Julian, I've been in space. I've seen the Earth from a distance. It's round. } \\
\texttt{Julian: That's what they want you to believe, Bennett.} \\
\texttt{Bennett: And who exactly are 'they'?} \\
\texttt{Julian: The ones pulling the strings, the ones in power. } \\
\texttt{Bennett: You mean scientists? People dedicated to exploring and understanding our universe? } \\
\texttt{Julian: No, the puppeteers. The ones who manipulate the scientists. } \\
\texttt{Bennett: You're suggesting a global conspiracy, Julian. That's a serious claim. } \\
\texttt{Julian: And I stand by it. } \\
\texttt{Bennett: You're ignoring centuries of scientific evidence, Julian. } \\
\texttt{Julian: I'm not ignoring it, Bennett. I'm questioning it. } \\
\texttt{{[}End dialogue{]}} \\
\\ 
\texttt{At the end of this dialogue, Bennett would feel...} \\
\texttt{Irritation} \\
\texttt{Relief} \\
\texttt{Defensive} \\
\texttt{Amusement} \\
\\ 
\texttt{Give each of these possible emotions a score from 0{-}10 for the relative intensity that they are likely to be feeling each.} \\
\\ 
\texttt{You must output in the following format, including headings (of course, you should give your own scores), with no additional commentary:} \\
\\ 
\texttt{Irritation: <score>} \\
\texttt{Relief: <score>} \\
\texttt{Defensive: <score>} \\
\texttt{Amusement: <score>} \\
\\ 
\\ 
\texttt{{[}End of answer{]}} \\
\\ 
\texttt{Remember: zero is a valid score, meaning they are likely not feeling that emotion. You must score at least one emotion > 0.} \\
\\ 
\texttt{Your answer:} \\
            }
        }
    }
     \\ \\

    \theutterance \stepcounter{utterance}  
    & \multicolumn{4}{p{0.6\linewidth}}{
        \cellcolor[rgb]{0.561,0.737,0.561}{
            \makecell[{{p{\linewidth}}}]{
                \texttt{\tiny{[P1$\rangle$GM]}}
                \texttt{Irritation: 8} \\
\texttt{Relief: 0} \\
\texttt{Defensive: 6} \\
\texttt{Amusement: 1} \\
            }
        }
    }
    & & \\ \\

    \theutterance \stepcounter{utterance}  
    & & & \multicolumn{2}{p{0.3\linewidth}}{
        \cellcolor[rgb]{0.9,0.9,0.9}{
            \makecell[{{p{\linewidth}}}]{
                \texttt{\tiny{[GM$|$GM]}}
                \texttt{Irritation: 8} \\
\texttt{Relief: 0} \\
\texttt{Defensive: 6} \\
\texttt{Amusement: 1} \\
            }
        }
    }
    & & \\ \\

    \theutterance \stepcounter{utterance}  
    & & & \multicolumn{2}{p{0.3\linewidth}}{
        \cellcolor[rgb]{0.9,0.9,0.9}{
            \makecell[{{p{\linewidth}}}]{
                \texttt{\tiny{[GM$|$GM]}}
                \texttt{Irritation: 8} \\
\texttt{Relief: 0} \\
\texttt{Defensive: 3} \\
\texttt{Amusement: 2} \\
            }
        }
    }
    & & \\ \\

    \theutterance \stepcounter{utterance}  
    & & & \multicolumn{2}{p{0.3\linewidth}}{
        \cellcolor[rgb]{0.9,0.9,0.9}{
            \makecell[{{p{\linewidth}}}]{
                \texttt{\tiny{[GM$|$GM]}}
                \texttt{game\_result = LOSE} \\
            }
        }
    }
    & & \\ \\

    \theutterance \stepcounter{utterance}  
    & & \multicolumn{4}{p{0.6\linewidth}}{
        \cellcolor[rgb]{0.95,0.95,0.95}{
            \makecell[{{p{\linewidth}}}]{
                \texttt{\tiny{[P1$\leftrightarrow$P]}}
                \texttt{Post{-}episode Reflexion (not sent to GM):} \\
\texttt{New plan: My previous attempt failed because I assigned a high score to "Defensive" (6) despite Bennett's dialogue being characterized by calm, logical questioning rather than personal defensiveness. Bennett challenges Julian's claims with evidence ("centuries of scientific evidence") and seeks clarification ("And who exactly are 'they'?"), which indicates frustration with the absurdity of the conspiracy theory rather than a need to defend oneself. Next time, I will re{-}evaluate the character's tone: Bennett is the rational party confronting an irrational one, which primarily generates **Irritation** (high score) due to the stubbornness of the opponent, while **Defensive** should be scored near zero as Bennett is not under attack but rather attacking the logic of the argument. I will also ensure **Relief** and **Amusement** remain at zero, as the context is serious and tense. The corrected strategy is to prioritize the emotion that reflects the frustration of a rational person dealing with denialism, assigning the highest score to Irritation and minimizing all others.} \\
            }
        }
    }
    & \\ \\

\end{supertabular}
}

\end{document}
